% Minimal style defining the commands a fine-screen resume template must use.
%
% The screener rewrites your resume by pattern, not by parsing LaTeX, so these
% four shapes are the contract:
%
%   \section{Title}          -- a section it may reorder, trim, or drop
%   \resumeSubheading{...}   -- one experience entry; must be followed by an
%                               item list and closed by \resumeItemListEnd
%   \resumeItemListStart/End -- the bullets under one entry, which it rewrites
%   \resumeBullet{...}       -- one bullet
%
% Restyle freely. Keep the command names and the nesting, or the rewriter will
% leave your resume untouched rather than corrupt it.

\usepackage[margin=1.6cm]{geometry}
\usepackage{enumitem}
\usepackage{titlesec}
\usepackage[hidelinks]{hyperref}

\pagestyle{empty}

\titleformat{\section}{\large\bfseries}{}{0pt}{}[\titlerule]
\titlespacing{\section}{0pt}{10pt}{6pt}

\newcommand{\contactSep}{\hspace{0.5em}\textbar\hspace{0.5em}}

% #1 role, #2 dates, #3 employer, #4 location
\newcommand{\resumeSubheading}[4]{%
  \vspace{4pt}%
  \noindent\textbf{#1} \hfill #2\\%
  \textit{#3} \hfill \textit{#4}%
  \vspace{2pt}%
}

\newcommand{\resumeBullet}[1]{\item\small{#1}}
\newcommand{\resumeItemListStart}{\begin{itemize}[leftmargin=*, itemsep=2pt, topsep=2pt]}
\newcommand{\resumeItemListEnd}{\end{itemize}}
\newcommand{\resumeSubHeadingListStart}{\begin{list}{}{\leftmargin=0pt}}
\newcommand{\resumeSubHeadingListEnd}{\end{list}}

\newcommand{\resumeHeadline}[1]{\begin{center}\small #1\end{center}\vspace{4pt}}
\newcommand{\resumeStatusLine}[1]{\begin{center}\small\textit{#1}\end{center}}
